\clearpage\SourceLeaf{11}\ArithmeticHead{4}{L}{BACHMANN, ÜBER GAUSS’ ZAHLENTHEORETISCHE ARBEITEN.}
\begin{center}\large\bfseries I. Die Disquisitiones arithmeticae.\end{center}
\begin{center}\textls[100]{a. Die Entstehung der D. A. Die Analysis Residuorum.}\end{center}
\begin{center}\bfseries 1.\end{center}
\noindent\hspace*{\parindent}Die D. A. erschienen im Sommer des Jahres 1801. In der Vorrede zu seinem Werke setzt \textsc{Gauss} den Beginn seiner Beschäftigung mit dessen Gegenstande in den Anfang des Jahres 1795. Dies ist nun freilich nicht ganz genau zu nehmen. Schon vorher hat \textsc{Gauss} sich viel mit rechnerischen Versuchen beschäftigt, die hauptsächlich auf die Teilbarkeit der ganzen Zahlen, insbesondere auf die Eigenschaften der Primzahlen, auf die Reste, welche die Potenzen andrer Zahlen durch solche geteilt lassen, u. dgl. m. gerichtet gewesen sind. \textsc{Gauss} selbst erwähnt in einem an \textsc{Encke} gerichteten Briefe (W. II, S. 444), dass er bereits (im Jahre 1792 oder 1793) als 15jähriger, als er sich die \textsc{Lambert}schen Supplemente zu den Logarithmentafeln verschafft habe, mit einer der verzwicktesten Aufgaben der Zahlentheorie, der Frequenz der Primzahlen, rechnerisch sich beschäftigt habe. Bald hernach schon hat \textsc{Gauss} diesem Briefe zufolge die Tafel solcher Frequenz, welche W. II, S. 435 abgedruckt ist, begonnen, während ihre vollständige Berechnung erst späterer Zeit angehört. Auch entstammt jener Zeit wohl schon die Tafel des quadratischen Charakters der Primzahlen (W. II, S. 399), jedenfalls aber die Tafel zur Verwandlung gemeiner Brüche in Dezimalbrüche (W. II, S. 411), die beide nur zum Teil den artt. 99 resp. 316 der D. A. angefügt sind. Von der letzteren sagt \textsc{Gauss} a. a. O.: quam integram sive etiam ulterius continuatam occasione data publici juris faciemus. Sie besteht aus zwei Teilen, deren zweiter in der Handschrift des Nachlasses den Vermerk trägt: explicitus October 11. 1795. Ob nun dieser zweite Teil die geplante weitere Fortsetzung ausmache oder nicht, jedenfalls gehört diese Tabelle dem genannten Jahre an; im ersteren Falle aber dürfte man schliessen, dass der bezügliche Teil der D. A. schon vor jenem Tage verfasst sein müsse. Dieser Ansicht zu Hilfe käme eine andere Tabelle mit dem Titel: »I. Ausschliessung gewisser Zahlen beim Aufsuchen der Faktoren«, welche in einem Hefte des Nachlasses mit der Aufschrift »Neue Beiträge zu den mathematischen Tafeln besonders zur Erleichte-\newpage
\SourceLeaf{12}\ArithmeticHead{5}{R}{DIE ENTSTEHUNG DER D. A. DIE ANALYSIS RESIDUORUM.}
\noindent rung der Zahlenrechnungen« und mit dem Datum März 1795 vorhanden ist. Sie bezieht sich auf die Methode der artt. 319—322 des sechsten Abschnitts der D. A., deren Inhalt hiernach, soweit er nicht die Theorie der quadratischen Formen voraussetzt, als von \textsc{Gauss} schon im Jahre 1795 gefunden angenommen werden darf. Die »Hülfstafel bei Auflösung der unbestimmten Gleichung \(A=fx^2+gy^2\) vermittelst der Ausschliessungsmethode« (W. II, S. 509), bei welcher schon der sechste Abschnitt der D. A. erwähnt wird, auf den sie sich bezieht, gehört erst späterer Zeit an.
